\section{Setting}\label{sec:setting}

\subsection{Nonlocal games, bipartite systems and types}

A nonlocal game consists of finite sets of questions $x$ for Alice and $y$ for Bob, finite
sets of answers $a$ and $b$, a probability distribution $\pi(x,y)$ and a rule
$V(a,b\,|\,x,y)\in\{0,1\}$. The referee draws $(x,y)$ from $\pi$; the players, who cannot
communicate, answer $a$ and $b$ and win if $V(a,b\,|\,x,y)=1$~\cite{CHTW04,BrunnerRMP}.

\emph{Finite systems.} Let $\psi\in\cH_A\otimes\cH_B$ be a unit vector. A strategy assigns to
each question a positive operator-valued measure (POVM), $\{A^x_a\}_a$ on $\cH_A$ and
$\{B^y_b\}_b$ on $\cH_B$. Its losing probability is $\ip{\psi}{L\psi}$ with the loss operator
\begin{equation}\label{eq:loss}
L=\sum_{x,y}\pi(x,y)\sum_{a,b\,:\,V(a,b|x,y)=0}A^x_a\otimes B^y_b,\qquad 0\le L\le 1 ,
\end{equation}
and the value of the game at $\psi$ is $\omega(\psi)=1-\inf\ip{\psi}{L\psi}$, the infimum being
over all strategies. Local ancillas in fixed pure states are included, since a measurement on
system plus ancilla is a POVM on the system.

\emph{Infinite systems.} The Hilbert space of an infinite system need not factorize. A bipartite
system is then a pair of von Neumann algebras $M_A$, $M_B$ on a separable Hilbert space $\cH$
that commute, $M_B\subseteq M_A'$, together with a unit vector
$\Omega\in\cH$. Haag duality is the stronger equality $M_B=M_A'$. Alice's POVMs are taken in $M_A$ and Bob's in $M_B$, the product
$A^x_a\otimes B^y_b$ in~\eqref{eq:loss} is replaced by $A^x_aB^y_b$, and $\omega(\Omega)$ is
defined as before. For $M_A=\mathcal{B}(\cH_A)\otimes 1$ and $M_B=M_A'$ this is the tensor-product definition. For a pure
state of a spin chain with quasi-local algebra $\fA=\fA_L\otimes\fA_R$ (left and right
half-chains) and GNS representation $(\pi,\cH,\Omega)$ we put $M_A=\pi(\fA_R)''$ and
$M_B=\pi(\fA_L)''$.

\emph{Types.} A von Neumann algebra with trivial centre is a factor. A factor has type I if it
is isomorphic to all bounded operators on a Hilbert space, type II if it is not of type I but
has a faithful normal semifinite trace, and type III otherwise. For a faithful normal state
$\varphi=\ip{\Omega}{\,\cdot\,\Omega}$ on $M$, use the cyclic GNS space $\overline{M\Omega}$,
where $\Omega$ is also separating. Tomita--Takesaki theory provides the modular operator
$\Delta_\varphi$ on this space and the modular group
$\sigma_t^\varphi(x)=\Delta_\varphi^{it}x\Delta_\varphi^{-it}$; the fixed points of
$\sigma^\varphi$ form the centralizer $M_\varphi=\{x:\varphi(xy)=\varphi(yx)\ \forall y\}$. For $M=\mathcal{B}(\cH_A)$ with a faithful density matrix of eigenvalues
$p_j>0$, the modular point spectrum consists of the ratios $p_i/p_j$; the
full spectrum is their closure. Connes's invariant $S(M)$ is the intersection
of modular spectra over faithful normal semifinite weights. It is $\{1\}$
for semifinite factors, $\{0,1\}$ for III$_0$,
$\{0\}\cup\lambda^{\Z}$ for III$_\lambda$, $0<\lambda<1$, and
$[0,\infty)$ for III$_1$~\cite{Connes73,TakesakiII}.

\subsection{The game}\label{sec:game}

The game $\GIII$ is the game called $G_{\mathrm{III}}$ in Ref.~\cite{TG}, and
Appendix~\ref{app:game} reproduces its definition. It is built from nine copies of the flagged
magic-square game of Man\v{c}inska~\cite{Mancinska26}, which extends the Mermin--Peres
magic-square game~\cite{Mermin90,Peres90} by one extra answer, called a flag, and two extra
questions per player. One further question to Bob, whose answer is a label in $\Z/9\Z$, links
the nine copies. Alice has 45 possible questions and Bob 46, each with at most nine answers,
and the 243 question pairs that occur are asked with equal probability.

If either local algebra admits a faithful normal semifinite trace, then every strategy
loses with positive probability~\cite[Theorem~1.1]{TG}. For a factor $M_A$ with separable
predual in standard form, the same theorem characterizes the existence of a normal state
with a perfect strategy on $(M_A,M_A')$ by an invariant normal state on the flow of weights.
Every type-III$_1$ factor meets that criterion. The robust lower comparison in
Theorem~\ref{thm:equiv} uses Ref.~\cite[Theorem~1.4]{TG}, with an explicit constant proved in
Appendix~\ref{sec:explicit-robust}.
For CHSH on a standard pair, van Luijk, Stottmeister, Werner and Wilming obtained
$\beta\ge 2\sqrt{2}-8\sqrt{\kappa(\varphi)}$, where $\beta$ is the CHSH correlation
coefficient and $\kappa$ is their worst-case marginal error in predual
norm~\cite[Theorem~J and Proposition~100]{vLSWW24}.

\begin{lemma}[Linked modular bridge]
\label{lem:linked-bridge}
Let $M_A,M_B$ be commuting von Neumann algebras on $\cH$ and $\Omega$ a unit vector
whose marginal $\varphi$ on $M_A$ is faithful. Suppose that POVMs in these algebras
win $\GIII$ with probability one at $\Omega$. For each $j\in\Z/9\Z$, let $F_j$ be
Alice's flag effect on $R_1^j$. The $F_j$ are orthogonal projections in the centralizer
of $\varphi$, with $\sum_jF_j=1$ and $\varphi(F_j)=1/9$. There are partial isometries
$w_j\in M_A$ such that, for all real $t$,
\begin{equation}\label{eq:bridge}
w_j^*w_j=F_j,\qquad w_jw_j^*=F_{j+1}+F_{j+2},\qquad \sigma^\varphi_t(w_j)=2^{it}w_j .
\end{equation}
\end{lemma}

\begin{proof}
Let $e$ project onto $\overline{M_A\Omega}$. This space reduces $M_A$, and the
representation there is faithful and normal. Compress Bob's effects to it by
$B\mapsto eBe|_{e\cH}$. They still form POVMs in the commutant, and every expectation
$\ip{\Omega}{AB\Omega}$ is preserved because $e\Omega=\Omega$ and $[e,A]=0$.
We work in this cyclic and separating representation of the faithful normal state
$\varphi$~\cite[Lemma~1.6]{TG}.

A perfect binary match of commuting effects $P,Q$ has zero expectation of
$(P-Q)^2+P(1-P)+Q(1-Q)$. Positivity gives $P\Omega=Q\Omega$, and faithfulness
makes Alice's effect $P$ a projection. Replacing $P\Omega$ by $Q\Omega$ and commuting
$Q$ past $X\in M_A$ gives $\varphi(PX)=\varphi(XP)$. The first label test therefore
gives centralizer projections $F_j$ with $F_j\Omega=B_j^{\mathrm{Lab}}\Omega$.
The label effects sum to one, so separation gives $\sum_jF_j=1$ and the $F_j$ are
orthogonal. The second label test matches Alice's $R_2^j$ outcome $+++$ to
$B_{j+1}^{\mathrm{Lab}}+B_{j+2}^{\mathrm{Lab}}$. On $\Omega$ this sum equals
$F_{j+1}+F_{j+2}$, so separation identifies Alice's outcome with that
projection~\cite[Section~2.4]{TG}.

Each copy is a perfect flagged magic-square strategy on this faithful normal support.
The checked square generates a centralizer copy of $M_4$ with unit $1-F_j$, and the four
nonflag outcomes of $R_2^j$ are rank-one projections in it~\cite[Lemma~2.1]{TG}.
The calibrated sign tests give a partial isometry $\beta_j$ from $F_j$ onto a rank-one
projection $q_j$ of that copy, with $\sigma_t^\varphi(\beta_j)=2^{it}\beta_j$~\cite[Lemma~2.2]{TG}.
Choose a partial isometry $v_j$ in that matrix algebra with initial projection $q_j$
and final projection $F_{j+1}+F_{j+2}$. Then $w_j=v_j\beta_j$ has the required initial
and final projections, and $\sigma_t^\varphi(w_j)=v_j\sigma_t^\varphi(\beta_j)$ gives the
last identity in~\eqref{eq:bridge}. The KMS identity gives
$2\varphi(F_j)=\varphi(F_{j+1})+\varphi(F_{j+2})$.
A largest weight forces both successors to be largest, so all nine weights are $1/9$.
\end{proof}

For a nonfaithful marginal with support $p\in M_A$, the bridge applies in the corner
$pM_Ap$ on $p\cH$. The equality $p\Omega=\Omega$ and commutation of $p$ with Bob's effects
preserve the correlations in this reduction. On the standard pair of a faithful marginal,
the maps $s_r=\sum_j w_j^*F_{j+r}$ for $r=1,2$ satisfy $s_r^*s_s=\delta_{rs}1$,
$\sum_rs_rs_r^*=1$ and $\sigma_t^\varphi(s_r)=2^{-it}s_r$. The local column isometries
with entries $s_r^*$ and $Js_r^*J$, where $J$ is the modular conjugation, extract a Bell
pair exactly. Indeed, Tomita's identity gives
$s_r^*(Js_s^*J)\Omega=2^{-1/2}s_r^*s_s\Omega=2^{-1/2}\delta_{rs}\Omega$.

\subsection{The embezzlement error}\label{sec:E}

For a unit vector $\psi\in\cH_A\otimes\cH_B$ the embezzlement error $\cE(\psi)$ is defined by
the infimum in~\eqref{eq:E}, taken over unitaries $U_A$ on $\cH_A\otimes\C^2$ and $U_B$ on
$\cH_B\otimes\C^2$; $\Phi_2=(\ket{00}+\ket{11})/\sqrt2$ is a Bell pair on the two added qubits.
This is embezzlement in the sense of van Dam and Hayden with $\psi$ as the catalyst, measured
by the squared distance; the maximal overlap, or root fidelity, is $1-\cE/2$. Let
$\lambda=(\lambda_j)_{j\ge1}$, $\lambda_j=\sqrt{p_j}$, be the decreasing sequence of Schmidt
coefficients of $\psi$, and define an isometry $D$ on $\ell^2(\N)$ by
\begin{equation}\label{eq:D}
(D\lambda)_{2j-1}=(D\lambda)_{2j}=\lambda_j/\sqrt2 .
\end{equation}

\begin{lemma}[Formula for the embezzlement error]\label{lem:formula}
$\cE(\psi)=\norm{\lambda-D\lambda}^2=2-\sqrt2\sum_{m\ge1}\sqrt{p_m\,p_{\lceil m/2\rceil}}$.
\end{lemma}

\begin{proof}
The Schmidt coefficients of $\psi\otimes\ket{00}$ are $\lambda$, and those of
$\psi\otimes\Phi_2$ are $D\lambda$, which is again decreasing. For two unit vectors with
decreasing Schmidt coefficients $\alpha$ and $\beta$, the supremum of the overlap
$\lvert\ip{\chi}{(U_A\otimes U_B)\chi'}\rvert$ over local unitaries is $\sum_j\alpha_j\beta_j$.
This is von Neumann's trace inequality~\cite{HornJohnson}; the bound is attained by aligning
the Schmidt bases in finite dimensions and approached in infinite dimensions. Hence the infimum
in~\eqref{eq:E} equals $2-2\ip{\lambda}{D\lambda}=\norm{\lambda-D\lambda}^2$, because $D$ is an
isometry.
\end{proof}


For a target of dimension $d$ the same argument gives
\begin{equation}\label{eq:Ed}
\cE_d(p)=\|(1-D_d)\lambda\|_2^2,
\qquad (D_d\lambda)_j=d^{-1/2}\lambda_{\lceil j/d\rceil}.
\end{equation}
The distinction between local-unitary and LOCC error matters for rates;
we use~\eqref{eq:Ed} throughout.
